\chapter{Backup and restore}
\label{ch:backup}

A self-hosted library has no cloud safety net. This chapter covers
what to back up, how often, and how to restore.

\section{What to back up}

Three artifacts together form a complete mybibli backup:

\begin{enumerate}
  \item \textbf{The MariaDB database.}\index{backup!database} Holds
        every title, volume, location, contributor, loan, user,
        audit log and setting. Lose this and you start from zero.
  \item \textbf{Cover images.}\index{backup!covers} Downloaded during
        cataloging and persisted in the Docker volume
        \texttt{mybibli\_covers} (or in your host bind-mount path if
        you switched per
        section~\ref{sec:install-covers-bind}). Lose this and the
        catalog still works, but every title page shows a generic
        placeholder until you re-fetch metadata title-by-title
        (issue~\#213 — pre-1.1.4 installs without this volume had
        covers disappear on every container upgrade).
  \item \textbf{The \texttt{.env} file.}\index{backup!.env} Holds
        passwords (database, API keys not yet migrated). Lose this
        and a re-deploy needs you to re-enter every credential.
\end{enumerate}

The Docker images themselves are \emph{not} backup-critical — you
can always pull them again from Docker Hub.

\section{Backup procedure}

\subsection*{Database — bundled MariaDB}

From the host running the compose stack:

\begin{lstlisting}
docker compose exec -T db mariadb-dump \
    -u root -p"$MYSQL_ROOT_PASSWORD" \
    --single-transaction --routines --triggers \
    mybibli \
    > backups/mybibli-$(date +%Y%m%d-%H%M%S).sql
\end{lstlisting}

The \texttt{--single-transaction} flag is important: it dumps a
consistent snapshot without locking the table for the duration of
the dump. \texttt{--routines --triggers} preserves any stored
routines (mybibli does not currently use any, but the flag is cheap
and future-proof).

Compress the dump if disk space is a concern:

\begin{lstlisting}
gzip backups/mybibli-*.sql
\end{lstlisting}

\subsection*{Database — external MariaDB}

Run \texttt{mariadb-dump} directly against the external server.
Synology DSM users can also schedule a \emph{Hyper Backup} task that
targets the MariaDB package — it produces a binary backup that
restores faster than a SQL dump.

\subsection*{Covers and \texttt{.env}}

For a host-mounted covers directory (bind-mount installs), a plain
\texttt{tar} archive is enough:

\begin{lstlisting}
tar czf backups/mybibli-files-$(date +%Y%m%d).tar.gz \
    .env "$COVERS_HOST_PATH"
\end{lstlisting}

For the default install (named Docker volume), back up the volume
through Docker — the host path of a named volume varies by Docker
version and is not stable, so don't try to \texttt{tar} it directly:

\begin{lstlisting}
docker run --rm \
    -v mybibli_covers:/source:ro \
    -v "$(pwd)/backups:/backup" \
    alpine \
    tar czf "/backup/mybibli-covers-$(date +%Y%m%d).tar.gz" -C /source .
tar czf "backups/mybibli-env-$(date +%Y%m%d).tar.gz" .env
\end{lstlisting}

The throwaway \texttt{alpine} container mounts the named volume
read-only, the host \texttt{backups/} directory writable, and writes
the tarball. Same shape works for restoring — see the Restore
section below.

\section{Recommended cadence}

\begin{itemize}
  \item \textbf{Daily} — automated database dump. Even on a
        seldom-changed catalog, a daily dump means worst-case data
        loss is bounded to one day.
  \item \textbf{Weekly} — covers and \texttt{.env}. These change
        only when you catalog new titles or rotate keys; weekly
        is plenty.
  \item \textbf{Off-site copy} — at least monthly. Rsync the
        \texttt{backups/} directory to a remote NAS, a USB drive
        you swap in / out, or a cloud-storage bucket. A fire that
        consumes the home server should not also consume your
        backups.
\end{itemize}

\begin{tip}
A backup that has never been restored is not a backup. Once a
quarter, restore your latest dump into a throwaway docker
environment and verify the title and loan counts match what
\texttt{/admin > Health} reports on production.
\end{tip}

\section{Restore procedure}

\subsection*{Database}

For the bundled-database setup:

\begin{lstlisting}
docker compose stop mybibli
docker compose exec -T db mariadb -u root -p"$MYSQL_ROOT_PASSWORD" \
    -e "DROP DATABASE mybibli; CREATE DATABASE mybibli CHARACTER SET utf8mb4;"
zcat backups/mybibli-20260513-1200.sql.gz | \
    docker compose exec -T db mariadb -u root -p"$MYSQL_ROOT_PASSWORD" mybibli
docker compose start mybibli
\end{lstlisting}

For external MariaDB, the same \texttt{mariadb} client commands work
against the external server.

\subsection*{Covers and \texttt{.env}}

For a host-bind-mount install:

\begin{lstlisting}
tar xzf backups/mybibli-files-20260513.tar.gz
\end{lstlisting}

For the default named-volume install, restore through a throwaway
container in the symmetric shape to the backup:

\begin{lstlisting}
docker compose stop mybibli
docker run --rm \
    -v mybibli_covers:/target \
    -v "$(pwd)/backups:/backup:ro" \
    alpine \
    sh -c "rm -rf /target/* && tar xzf /backup/mybibli-covers-20260513.tar.gz -C /target"
tar xzf backups/mybibli-env-20260513.tar.gz
docker compose start mybibli
\end{lstlisting}

Restart the stack so the mybibli container picks up the restored
\texttt{.env}.

\section{Migration to a new host}

The backup procedure is also the migration procedure:

\begin{enumerate}
  \item On the old host: produce a fresh database dump and tar
        archive.
  \item Transfer both to the new host.
  \item On the new host: install Docker, pull the compose file from
        the release matching the dump's version, place \texttt{.env}
        and \texttt{covers/} alongside, restore the database, start
        the stack.
  \item Verify by counting titles on \texttt{/admin > Health} and
        spot-checking a few title pages.
\end{enumerate}

\begin{warning}
The dump is tied to the \emph{schema} version. If the source host
ran 1.1.0 and the new host pulls 1.5.0, the new host's migrations
will upgrade the restored database in place — that is fine. The
reverse (restoring a 1.5.0 dump onto a 1.1.0 install) is \emph{not}
supported and will fail at boot. Always restore on a host running
the same or newer mybibli release.
\end{warning}
